% =====================================================================
%  config/preambule.tex
%  Réglages généraux : classe, géométrie, paquets de base.
%  Ce fichier ne charge ni les polices, ni les langues, ni le style
%  visuel FTI : chacun de ces aspects a son propre fichier de config.
% =====================================================================

% ---------------------------------------------------------------------
%  Vérification du moteur : ce modèle est prévu pour XeLaTeX.
%  (fontspec + polyglossia + rendu Unicode natif de l'arabe, du CJK…)
% ---------------------------------------------------------------------
\usepackage{iftex}
\RequireXeTeX % stoppe la compilation avec un message clair si ce n'est pas XeTeX

% ---------------------------------------------------------------------
%  Encodage d'entrée : avec XeLaTeX tout est UTF-8 nativement.
%  (Pas de \usepackage[utf8]{inputenc} ni de fontenc : inutiles/nuisibles.)
% ---------------------------------------------------------------------

% ---------------------------------------------------------------------
%  Géométrie de la page
%  typearea est déjà chargé par scrbook ; on ne le recharge pas (cela
%  provoquerait un « option clash »). On règle l'empagement via KOMAoptions.
%    BCOR = correction de reliure (marge intérieure supplémentaire).
%    DIV  = finesse de l'empagement KOMA (plus grand = marges + petites).
% ---------------------------------------------------------------------
\KOMAoptions{DIV=12,BCOR=8mm}
\recalctypearea

% ---------------------------------------------------------------------
%  Paquets de base
% ---------------------------------------------------------------------
\usepackage{xcolor}            % couleurs (utilisé par style-fti.tex)
\usepackage{graphicx}          % \includegraphics (logo, figures)
% microtype améliore la micro-typographie mais est facultatif : on ne
% le charge que s'il est installé (paquet Fedora : texlive-microtype).
\IfFileExists{microtype.sty}{\usepackage{microtype}}{}
\usepackage{booktabs}          % \toprule \midrule \bottomrule
\usepackage{tabularx}          % colonnes de largeur automatique
\usepackage{ragged2e}          % \RaggedRight, justification fine
\usepackage{setspace}          % \onehalfspacing pour le corps de la thèse
\usepackage{enumitem}          % listes personnalisables
\usepackage{csquotes}          % \enquote{} — guillemets adaptés à la langue

% ---------------------------------------------------------------------
%  Compatibilité TeX Live 2026 (arabe dans les tableaux)
%  Le noyau LaTeX 2026 a retiré la « fake math » autour des structures
%  tabulaires. Or le paquet « bidi » (chargé par polyglossia pour l'arabe,
%  voir config/langues.tex) appelle encore \UseMathForPositioningText,
%  désormais indéfini : dès qu'un tableau (tabular/tabularx) apparaît, on
%  obtient « Undefined control sequence » — y compris sur Overleaf (TL2026).
%  Ce shim rétablit la macro. Il est défini AVANT le chargement de bidi et
%  reste inoffensif sur les versions antérieures de TeX Live.
%  Réf. : https://www.overleaf.com/blog/tex-live-2026-is-now-available
%         (section « polyglossia, xepersian, and array »).
\makeatletter
\providecommand\UseMathForPositioningText{\m@th}
\makeatother

% Interligne : 1,5 est l'usage courant pour une thèse.
\onehalfspacing

% ---------------------------------------------------------------------
%  Coupures de ligne et débordements (« Overfull \hbox »)
%  ---------------------------------------------------------------------
%  Les noms d'auteurs en petites capitales issus de biblatex (p. ex.
%  « Schlechtweg ») et les liens hyperref ne se coupent pas toujours en
%  fin de ligne, ce qui peut faire déborder le texte dans la marge.
%  - \emergencystretch autorise TeX à relâcher l'espacement en dernier
%    recours plutôt que de déborder ;
%  - on autorise la césure à l'intérieur des liens hyperref (citations).
\emergencystretch=3em
\hfuzz=1pt
\hbadness=1500
% Autorise le retour à la ligne dans les hyperliens (citations, URL).
\usepackage{etoolbox}
\appto\UrlBreaks{\do\-}


% ---------------------------------------------------------------------
%  Titres et table des matières (réglages KOMA)
% ---------------------------------------------------------------------
\setcounter{tocdepth}{2}       % profondeur de la table des matières
\setcounter{secnumdepth}{3}    % profondeur de la numérotation des titres

% ---------------------------------------------------------------------
%  Liens hypertextes et métadonnées PDF.
%  hyperref doit être chargé tard ; polyglossia/bidi sont chargés
%  dans langues.tex, avant l'appel à ce bloc dans these.tex.
% ---------------------------------------------------------------------
\usepackage[
  unicode=true,                % métadonnées PDF en Unicode (titres non latins OK)
  hidelinks,                   % pas de cadres colorés autour des liens
  pdfusetitle,                 % titre/auteur PDF depuis \title/\author
  bookmarksnumbered=true,
  breaklinks=true,             % autorise un lien (citation, URL) à passer à la ligne
]{hyperref}

% Signets PDF corrects même avec des caractères Unicode.
\usepackage{bookmark}
